\documentclass[12pt]{article}

\usepackage{fontawesome}
\usepackage{hyperref}
\usepackage{xurl}
\usepackage[tagged, highstructure]{accessibility}
\usepackage{bookmark}
\usepackage{enumitem}

\hypersetup{
    colorlinks=false,
    pdfborder={0 0 0},
    pdftitle={Lab Template},
    pdfauthor={Adrianna Holden-Gouveia},
    pdfsubject={Linux Administration - Lab Assignment},
    pdflang={en-US},
    pdfkeywords={lab assignment, template}
}

% Configure PDF bookmarks for navigation
\bookmarksetup{
    numbered,
    open,
}

% Configure list spacing for better accessibility
\setlist{nosep}



\title{Lab Template}
\author{
        Adrianna Holden-Gouveia \\
        Website: \href{https://aholdengouveia.name}{aholdengouveia.name}\\ 
        \date{\vspace{-5ex}}
        %Email: \href{mailto:admin@aholdengouveia.name}{admin@aholdengouveia.name} \\
%        \faLinkedin{: aholdengouveia} \\
        \faGithub {: aholdengouveia} \\
       % \faTwitter {: aholdengouveia} \\
        }

%S\date{\today}


\begin{document}    

\maketitle

%\begin{abstract}
%This is a template for Linux Administration Lab work
%\end{abstract}
%\tableofcontents


\section*{Accessibility Notice}
This document is also available in HTML format at:

Link: \href{https://aholdengouveia.name/LinuxAdmin/labexcercises/labtemplate.html}{\textbf{\nolinkurl{https://aholdengouveia.name/LinuxAdmin/labexcercises/labtemplate.html}}}

The HTML version provides enhanced accessibility features including keyboard navigation, screen reader support, responsive design, dark mode support, and high contrast options.

\section*{Objectives:}
\begin{itemize}
    \item The objective of this lab is to ... Students will ... These skills are ...
\end{itemize}

% Background: 2 to 3 sentences of Acme Corp context. Why does the team need this work?
% If the lab uses the servers, name them the same way as the other labs.
\section*{Background}
Acme Corp ... You will be working on the two servers you set up in the Server Setup lab: your Red Hat based server and your second server.

% Requirements: the actual work. Keep them as open as you want students to handle on their own.
\section*{Requirements}
\begin{enumerate}
    \item Requirement 1
    \item Requirement 2
    \item Requirement 3
\end{enumerate}

% Acceptance Criteria: how we know the work is done (proof, testing, citations).
% Keep ONE of the two citation rules below and delete the other.
\section*{Acceptance Criteria}
\begin{itemize}
    % Full rule, for labs with real scripts:
    \item \textbf{Everything must be cited.} Every command, statement, and line of code in your scripts and documents must be cited, either with a link to a reference or with a note naming the previous class where you learned it, for example: \texttt{\# Learned in CIS 117 Intro to Linux}. In scripts, put the citation in a comment next to the code. Anything without a citation will not count.
    % Lighter rule, for labs of single commands (grep, awk, sed):
    %\item If you had to look anything up to figure out a command, cite your source with a link to the reference.
    \item Everything is tested on both servers, with screenshots as proof.
    \item Everything is in your own words.
    \item Every terminal screenshot shows your prompt with your username and hostname, and either the output of \texttt{date} or a note beside the terminal with your name and the term. Hostnames must match your server inventory from the Server Setup lab.
    \item Your GitHub repository shows your work over time: several commits across more than one day, each with a message describing the change, not a single upload of finished code.
\end{itemize}

% Deliverables: what to hand in. Unnumbered, like industry.
\section*{Deliverables}
Audience: ... Assume technical knowledge but not expertise.  Screenshots are helpful to go with your descriptions.

Make sure your deliverables include documents for BOTH servers, clearly labeled with what the document is and which server it's for.
\begin{itemize}
    \item Deliverable 1
    \item Deliverable 2
    \item Scripts, well commented and tested, with a link to your GitHub where you've uploaded them.
\end{itemize}

% Resources: your website, tutorials, and any linked files (full links).
\section*{Resources}
References, a video, a PowerPoint and some notes are available at my website
Link: \href{https://www.aholdengouveia.name/LinuxAdmin/}{\textbf{\nolinkurl{https://www.aholdengouveia.name/LinuxAdmin/}}}

% ============================================================
% EXAMPLE: Accessible Table (if needed)
% Uncomment and modify the lines below to add a table
% This ensures proper screen reader support
% ============================================================
%\begin{table}[h]
%\caption{Command Options}
%\begin{tabular}{|l|l|}
%\hline
%\textbf{Option} & \textbf{Description} \\
%\hline
%-F & Field separator \\
%\hline
%-v & Variable assignment \\
%\hline
%\end{tabular}
%\end{table}

\end{document}